\section{Introduction}\label{sec:intro}

Empirical machine learning advances one comparison at a time. A new gradient-boosting configuration is reported to beat a random forest, a deep architecture to beat gradient boosting, one feature set to beat another. Such claims rest on a small number of summary statistics (an average accuracy, an F1-score, a root mean squared error) computed by cross-validation, and, in the better cases, on a statistical test or a confidence interval for the difference. The test is meant to answer a simple question: \emph{is the observed improvement larger than what the randomness of the data and of the validation procedure could produce on its own?}

The standard tools for that question~\cite{dietterich1998,nadeau2003,demsar2006,benavoli2017} were developed for observations that are independent and identically distributed (i.i.d.). Many of the data sets to which machine learning is now applied are not. Meteorological, environmental, agricultural, epidemiological and mobility data are collected at locations and times that are close to one another, and nearby observations share unmeasured drivers. Two consequences are well documented. First, random $K$-fold cross-validation places near-duplicates of test observations in the training set, so that it measures interpolation rather than prediction at new locations or times, and it is optimistic when the latter is the target~\cite{roberts2017,ploton2020,meyer2018,bergmeir2012}. Second, the \emph{errors} of a fitted model are themselves dependent, so that $n$ test observations carry much less information than $n$ independent ones; this second problem has been addressed for fixed forecasts and spatial predictions~\cite{diebold1995,hering2011}, but rarely for learning procedures evaluated by cross-validation.

The first problem has a well-known remedy: blocked (spatial, temporal or spatiotemporal) cross-validation~\cite{roberts2017,valavi2019,linnenbrink2024}. The second problem is the subject of this paper. Blocking changes \emph{what} is estimated; it does not tell us \emph{how precisely} it is estimated. A practitioner who computes a blocked cross-validation score for two models and then applies a paired $t$-test, a bootstrap that resamples observations, or a test on fold-wise scores still relies on independence assumptions that the data violate. Whether an apparent improvement is ``statistically real'' then depends as much on these hidden assumptions as on the models.

Our starting point is the per-observation out-of-fold \emph{loss differential} $d_i = \ell(y_i,\hat f_A(x_i)) - \ell(y_i,\hat f_B(x_i))$ between two models $A$ and $B$. The mean of the $d_i$ estimates the expected gain of $A$ over $B$, so that the comparison becomes a problem of inference on the mean of a dependent random field. Two sources of dependence must be handled at once. (i)~Test-point dependence: the $d_i$ inherit the spatial and temporal autocorrelation of the residuals of both models. (ii)~Training dependence: all observations of a fold are predicted by the same fitted models, and the models of different folds are fitted on overlapping training sets. Existing procedures address at most one of these sources. Heteroskedasticity- and autocorrelation-consistent (HAC) variance estimators, as in the Diebold--Mariano test for forecasts~\cite{diebold1995}, and the parametric spatial prediction comparison test of Hering and Genton~\cite{hering2011} address (i) for fixed predictors, whereas the corrected resampled $t$-test of Nadeau and Bengio~\cite{nadeau2003} addresses (ii), but only for i.i.d. data.

\paragraph{Contributions.} The paper makes five contributions.
\begin{enumerate}
\item We formalise model comparison under spatiotemporal dependence as inference on the mean of a loss-differential field with correlated fold-level training effects. We show (Proposition~\ref{prop:pairedt}) that the paired $t$-test on fold scores underestimates the variance by exactly the average covariance between fold means, and we derive (Proposition~\ref{prop:bias}) the exact bias of a variance estimator combining test-sampling and training components, which identifies the conditions under which it is unbiased and the direction of its error otherwise.
\item We propose \emph{DA-P}, a dependence-aware test for two models evaluated by blocked cross-validation. It combines (a) a test-sampling variance computed from a separable correlation model fitted to the variograms of the fold-centred loss differentials; (b) a method-of-moments estimate of the training-induced fold variance, with an explicit assumption on the correlation of training effects; and (c) a Welch--Satterthwaite reference distribution. All steps use a single cross-validation run.
\item We attach to the test a three-way decision rule that declares a comparison \emph{indeterminate} when its significance depends on the unverifiable correlation of training effects, and we evaluate it, together with a range-based diagnostic, by simulation.
\item We assess conventional and dependence-aware procedures (paired $t$, i.i.d. bootstrap, corrected resampled $t$, block bootstrap, linear mixed model, Diebold--Mariano--HAC, a space--time version of the Hering--Genton test, a kernel variant DA-K, and DA-P) in simulation studies in which the truth is known exactly. These studies cover the sensitivity to the assumed training correlation, the robustness to anisotropy, covariance shape, heteroscedasticity, non-stationarity and irregular sampling, and gradient-boosting learners whose expected performances are equal by an exchangeability argument.
\item We apply all procedures to three public data sets (hourly weather in New York, hourly pedestrian counts in Melbourne and yield-monitor maize data in Argentina) under random, blocked and forward-chaining validation, with a correction for multiple comparisons and a ``seed-null'' diagnostic, and we derive practical recommendations. An open-source Python implementation reproduces all results.
\end{enumerate}

The emphasis of the paper is thus not on whether spatial or temporal validation is preferable, a question already addressed in the literature, but on the \emph{validity of the inference} drawn from a given validation design. In the terms of the call for this special issue, it contributes to measuring the trustworthiness of comparisons between machine learning models in real applications.

The rest of the paper is organised as follows. Section~\ref{sec:related} reviews related work. Section~\ref{sec:problem} sets up the problem and analyses conventional tests. Section~\ref{sec:method} presents DA-P and its decision rule. Sections~\ref{sec:sim1} and~\ref{sec:sim2} report the simulation studies, Section~\ref{sec:real} the case studies, and Section~\ref{sec:discussion} discusses recommendations and limitations.
